% swim-times-sdd.tex
% Software Design Document for the swim-times program
% IEEE Std 1016-2009 structure
% Compiled with: pdflatex swim-times-sdd.tex

\documentclass[12pt]{article}

\usepackage{geometry}
\geometry{letterpaper, margin=1in}

\usepackage{hyperref}
\usepackage{booktabs}
\usepackage{array}
\usepackage{longtable}
\usepackage{enumitem}
\usepackage{titlesec}
\usepackage{fancyhdr}
\usepackage{xcolor}
\usepackage{listings}

\definecolor{castblue}{RGB}{0,48,135}
\definecolor{lightgray}{gray}{0.95}

\lstset{
  basicstyle=\ttfamily\small,
  backgroundcolor=\color{lightgray},
  frame=single,
  framesep=4pt,
  rulecolor=\color{lightgray},
  breaklines=true,
  showstringspaces=false,
  language=C,
  keywordstyle=\bfseries\color{castblue},
  commentstyle=\itshape\color{gray},
}

\pagestyle{fancy}
\fancyhf{}
\renewcommand{\headrulewidth}{0.4pt}
\fancyhead[L]{\small College Area Swim Team}
\fancyhead[R]{\small SDD: swim-times}
\fancyfoot[C]{\thepage}

\titleformat{\section}{\large\bfseries\color{castblue}}{\thesection}{0.5em}{}[\titlerule]
\titleformat{\subsection}{\normalsize\bfseries}{\thesubsection}{0.5em}{}

\newcommand{\code}[1]{\texttt{#1}}
\newcommand{\req}[1]{\textbf{#1}}

\begin{document}

% -----------------------------------------------------------------
% Title block
% -----------------------------------------------------------------
\begin{titlepage}
  \centering
  \vspace*{2cm}
  {\Huge\bfseries\color{castblue} Software Design Document\\[0.5em]}
  {\Large swim-times: Automated Swimmer-Times Retrieval System}\\[2em]
  {\large College Area Swim Team (CAST)}\\[0.5em]
  {\normalsize Prepared for: CAST Coaching Staff and Maintainers}\\[0.5em]
  {\normalsize Prepared by: CAST Technology Working Group}\\[2em]
  \begin{tabular}{ll}
    Document identifier: & CAST-SDD-001                                \\
    Revision:            & 2.0                                         \\
    Date:                & July 11, 2026                               \\
    Status:              & Released                                    \\
    Governing standard:  & IEEE Std 1016-2009 (Software Design Descriptions) \\
  \end{tabular}
  \vfill
  {\small\textit{This document is prepared in accordance with the principles
  of rational documentation articulated in D.\,L.\,Parnas and P.\,C.\,Clements,
  ``A Rational Design Process: How and Why to Fake It,''
  \textit{IEEE Transactions on Software Engineering}, vol.\,SE-12, no.\,2,
  pp.\,251--257, February 1986.}}
\end{titlepage}

\tableofcontents
\clearpage

% =================================================================
\section{Introduction}
% =================================================================

\subsection{Purpose}

This Software Design Document (SDD) describes the architecture
and detailed design of the \textbf{swim-times} program.  It
satisfies the design-description function of IEEE~Std~1016-2009
by presenting a set of design views, each tailored to a specific
stakeholder concern.  Together with the Concept of Operations
(CAST-CONOPS-001) and the Software Requirements Specification
(CAST-SRS-001), this document forms the complete development
baseline for swim-times.

\subsection{Scope}

This SDD covers all code that is tangled out of the literate
program \code{swim2/swim-times.w}.  It does not cover the
build system, deployment, or the companion Pandoc-based PDF
report (\code{swim/swim-times.md}); those are described
elsewhere.

\subsection{Relationship to the Parnas--Clements model}

Per Parnas and Clements (1986), this document presents the
design as if it had been derived top-down from requirements,
even though in practice the program evolved through exploratory
prototyping against an undocumented third-party API.  Each
design decision is presented in the form a reader would find
most useful, with the rationale recorded so that future changes
can be evaluated against the original intent.

% =================================================================
\section{Design Stakeholders and Concerns}
% =================================================================

\begin{center}
\begin{longtable}{p{3.5cm}p{3.5cm}p{6cm}}
  \toprule
  \textbf{Stakeholder} & \textbf{Concern} & \textbf{Addressed by view} \\
  \midrule\endhead
  Maintainers     & Module decomposition  & Composition view (\S4.2)        \\
  Maintainers     & Data structures       & Information view (\S4.4)        \\
  Maintainers     & Algorithmic detail    & Algorithm view (\S4.7)          \\
  Coaches/parents & Behaviour at runtime  & Interaction view (\S4.6)        \\
  Integrators     & External interfaces   & Context view (\S4.1) and Interface view (\S4.5) \\
  Auditors        & Security boundaries   & Resource view (\S4.8)           \\
  Reviewers       & Design rationale      & Design rationale (\S5)          \\
  \bottomrule
\end{longtable}
\end{center}

% =================================================================
\section{Design Approach}
% =================================================================

\subsection{Architectural style}

swim-times is a single-process, single-threaded
\emph{pipes-and-filters} architecture realised within one address
space: command-line options flow into an option-parsing stage,
which feeds a swimmer-resolution stage, which feeds a per-event
fetch-parse-sort stage, whose output is serialised by a
formatter to standard output.  Each stage is a CWEB chunk
(or small group of chunks) with a single responsibility.

\subsection{Key design decisions}

\begin{description}[style=nextline]
  \item[D-1: CWEB literate source.]
        The canonical source is \code{swim-times.w};
        \code{swim-times.c} is a build artefact.  This satisfies
        \req{REQ-D-02} and \req{REQ-A-04}.

  \item[D-2: 24-line chunk discipline.]
        No CWEB chunk exceeds twenty-four lines.  Larger logical
        units are decomposed into named sub-chunks.  This satisfies
        \req{REQ-D-03} and improves readability for the typeset
        documentation.

  \item[D-3: String-scanning JSON parser.]
        A purpose-built scanner (\code{scan\_string}) is used in
        place of a third-party JSON library.  This satisfies
        \req{REQ-D-04}.  The same generic ``find key, take the
        following value'' scanner served the old Sisense JAQL
        responses and now serves the named-field REST responses
        unchanged (see D-6).  Every field the program consumes ---
        meet name, event code, date, time, and standard --- is a JSON
        string, so a single string extractor suffices; the former
        numeric \code{scan\_long} helper was retired with the GET
        best-times feed it served.

  \item[D-4: Dynamic anonymous credentials (no bearer token).]
        The program no longer embeds a bearer token.  The USA
        Swimming times service authenticates anonymous callers
        with three headers --- \code{Usas-Sub-Id: Anonymous},
        \code{AppName: DataHub}, and a client-minted
        \code{Device-Id}.  \code{device\_id()} generates the
        \code{Device-Id} at run time in the format the server
        validates, so the ``token'' is obtained dynamically on
        every run.  See \S\ref{sec:why-anon}.

  \item[D-5: Insertion sort.]
        Per-event row counts rarely exceed a few dozen, so
        $O(n^2)$ insertion sort is used in preference to a more
        complex algorithm.  Simpler code, fewer failure modes.

  \item[D-6: REST times API, not Sisense JAQL.]
        The data hub retired the Sisense analytics endpoint this
        program originally targeted.  The fetch path was migrated
        to the public \code{times-api.usaswimming.org} REST
        service: \code{GetMembersForFilters} resolves a name to a
        \code{memberId}, and the \code{POST /BestTimes} query
        returns that member's best time per course for each
        stroke-and-distance pair, complete with swim date, meet
        name, and motivational standard.  See
        \S\ref{sec:why-rest} and \S\ref{sec:why-bestonly}.

  \item[D-8: Rich best-times query for date, meet, and standard.]
        The lightweight \code{GET /GetBestTimesForMember} feed used
        earlier returned only stroke, distance, course, and time.
        To satisfy \req{REQ-F-32} and \req{REQ-F-50} (report each
        time's date, meet, and B/BB/A/AA/AAA/AAAA standard) the fetch
        path was moved to
        the heavier \code{POST /BestTimes} query, which is still
        anonymous but carries \code{swimDate}, \code{meetName}, and
        \code{timeStandard}.  One POST is issued per distinct
        stroke-and-distance pair and the results are cached per
        member.  See \S\ref{sec:why-bestonly}.

  \item[D-7: Id-based swimmer selection.]
        Swimmers are selected by a unique string \code{id} carried
        on each \code{SWIMMERS} entry rather than by a per-swimmer
        option bit, so the roster can grow past the width of the
        flag word.  See \S\ref{sec:why-ids}.
\end{description}

% =================================================================
\section{Design Views}
% =================================================================

% --------------------------- 4.1 ---------------------------------
\subsection{Context view}

\paragraph{External actors.}
\begin{itemize}
  \item \textbf{Operator} --- a coach, parent, or maintainer
        invoking the binary from a shell.
  \item \textbf{USA~Swimming times API} --- the upstream HTTPS
        service at \texttt{times-api.usaswimming.org}.
  \item \textbf{Standard output / standard error} --- the
        downstream sinks that receive table or CSV output and
        diagnostics, respectively.
\end{itemize}

\paragraph{Trust boundary.}
The only trusted external system is the times API, which
serves only public competitive records.  No user-supplied data
crosses a trust boundary; command-line arguments are consumed
internally.

% --------------------------- 4.2 ---------------------------------
\subsection{Composition view}

The program is decomposed into the following top-level CWEB
chunks (the order in which they appear in the tangled output):

\begin{center}
\begin{longtable}{ll}
  \toprule
  \textbf{Chunk} & \textbf{Responsibility} \\
  \midrule\endhead
  \code{Includes}                 & Standard headers, \code{<curl/curl.h>}, \code{<libpq-fe.h>} \\
  \code{Global constants}         & Numeric limits, \code{TIMES\_API} base URL, event table \\
  \code{Option flags}             & \code{OPT\_*} bitmasks and global option state \\
  \code{Type definitions}         & \code{Buffer}, \code{TimeRow}, \code{Swimmer}
                                    structs                                          \\
  \code{HTTP utilities}           & \code{write\_cb}, \code{base64}, \code{device\_id}, \code{http\_request} \\
  \code{JSON scanner}             & \code{scan\_string}                            \\
  \code{Database connection}      & \code{db\_conn\_string}, \code{db\_open},
                                    \code{db\_close}, \code{db\_insert\_row}        \\
  \code{Person lookup}            & \code{lookup\_member\_id}                      \\
  \code{Times fetch}              & \code{insertion\_sort}, \code{time\_to\_seconds},
                                    \code{normalize\_date}, \code{cache\_add},
                                    \code{fetch\_pair}, \code{build\_member\_cache},
                                    \code{fetch\_times}                            \\
  \code{Offline fetch}            & \code{offline\_fetch} (libpq \code{SELECT})    \\
  \code{Main function}            & Option parsing, swimmer table, \code{main}     \\
  \bottomrule
\end{longtable}
\end{center}

Each top-level chunk is itself a composition of named sub-chunks
to honour the 24-line discipline (D-2).  For example,
\code{http\_request} expands to \code{Initialize curl handle} +
\code{Issue HTTP request and return}; \code{fetch\_times} rebuilds the
per-member cache (via \code{build\_member\_cache} $\to$
\code{fetch\_pair} $\to$ \code{Parse BestTimes records}) when the member
changes, then selects the cached row and hands it to the shared
sort/emit sub-chunks.

% --------------------------- 4.3 ---------------------------------
\subsection{Dependency view}

\begin{center}
\begin{tabular}{ll}
  \toprule
  \textbf{Caller}       & \textbf{Callee}                              \\
  \midrule
  \code{main}           & \code{parse\_opts\_str}, \code{parse\_events\_str}, \code{print\_usage} \\
  \code{main}           & \code{db\_open}, \code{db\_close} (under \code{store}) \\
  \code{main}           & \code{lookup\_member\_id}, \code{fetch\_times} (online) \\
  \code{main}           & \code{offline\_fetch} (under \code{offline})            \\
  \code{lookup\_member\_id} & \code{http\_request}, \code{scan\_string}          \\
  \code{fetch\_times}   & \code{build\_member\_cache}, \code{insertion\_sort}      \\
  \code{build\_member\_cache} & \code{fetch\_pair} $\to$ \code{http\_request},
                          \code{scan\_string}, \code{cache\_add}                 \\
  \code{cache\_add}     & \code{time\_to\_seconds}, \code{normalize\_date}         \\
  \code{offline\_fetch} & \code{PQexecParams}, \code{PQgetvalue},
                          \code{insertion\_sort}                          \\
  \code{db\_open}/\code{db\_close} & \code{PQconnectdb}, \code{PQprepare},
                          \code{PQfinish} (libpq)                         \\
  \code{db\_insert\_row} & \code{PQexecPrepared} (libpq)                  \\
  \code{http\_request}  & \code{libcurl} easy interface, \code{write\_cb} \\
  \bottomrule
\end{tabular}
\end{center}

The dependency graph is acyclic.  The two linked external
dependencies are \code{libcurl} (HTTP) and \code{libpq} (Postgres
client library), both linked at build time.  Postgres is reached
\emph{programmatically} via libpq calls --- no external
\texttt{psql} process or shell wrapper is involved (this
realises requirement~\#3 of \texttt{requirements4.txt}).

% --------------------------- 4.4 ---------------------------------
\subsection{Information view (data structures)}

\paragraph{\code{Buffer}.}
A growable heap buffer used to accumulate libcurl response chunks.
Always null-terminated so it can be treated as a C string.

\begin{lstlisting}
typedef struct {
    char  *data;
    size_t size;
} Buffer;
\end{lstlisting}

\paragraph{\code{TimeRow}.}
One swim-time record.  \code{sort\_key} is the swim time in seconds
(\code{time\_to\_seconds}) used to order rows fastest-first; the
remaining fields are fixed-size character arrays.  The
\code{POST /BestTimes} feed (\S\ref{sec:why-bestonly}) populates every
field; \code{standard} is left empty only when a swim earned no
motivational standard.

\begin{lstlisting}
typedef struct {
    double sort_key;
    char   time[32];
    char   date[16];
    char   standard[48];
    char   meet[256];
} TimeRow;
\end{lstlisting}

\paragraph{\code{Swimmer}.}
A static configuration record describing how to identify each
tracked swimmer.  \code{id} is the unique command-line selection
token (\S\ref{sec:why-ids}); \code{search\_query} and
\code{match\_substr} resolve the member via \code{GetMembersForFilters};
\code{date\_min}/\code{date\_max} are an optional age window, retained
for signature compatibility but not applied online
(\S\ref{sec:why-bestonly}).  The same physical swimmer may appear in
\code{SWIMMERS[]} more than once with distinct ids --- this is how the
three Ledecky age-group views are declared.

\begin{lstlisting}
typedef struct {
    const char *id;            /* unique selection id, e.g. "stella" */
    const char *search_query;  /* member-search "name" value         */
    const char *match_substr;  /* lowercase substring filter         */
    const char *date_min;      /* "YYYY-MM-DD" inclusive (or NULL)    */
    const char *date_max;      /* "YYYY-MM-DD" inclusive (or NULL)    */
} Swimmer;
\end{lstlisting}

\paragraph{Database schema (\texttt{schema.sql}).}
Persisted rows live in the \texttt{swim\_times} table, whose
columns mirror the seven fields written by the DB-pipe emitter
plus an identity primary key and a \texttt{fetched\_at}
timestamp.  A UNIQUE constraint on
\texttt{(swimmer, event, swim\_time, swim\_date, meet)} prevents
duplicate inserts on repeated \code{store} runs (REQ-F-73).
Two indexes (\texttt{(swimmer, event)} and
\texttt{(event, sort\_key)}) accelerate the \code{offline}
\texttt{SELECT} (REQ-F-71) which filters on swimmer name and
event code and orders by \texttt{sort\_key}.

\paragraph{Global state.}
Module-level variables hold parsed command-line state and the
best-times cache:

\begin{itemize}
  \item \code{g\_opts} --- bitmask of behaviour \code{OPT\_*} flags
        (\code{fastest}, \code{csv}, \code{store}, \code{offline}).
  \item \code{g\_sel[MAX\_SEL]}, \code{g\_nsel} --- selected swimmer
        ids; empty means ``all swimmers''.
  \item \code{g\_events[NUM\_EVENTS][32]}, \code{g\_nevents} ---
        user-requested event codes and their count.
  \item \code{g\_cache\_member}, \code{g\_cache[]}, \code{g\_cache\_n}
        --- the per-member best-times cache: one \code{TimeRow} per
        event, filled by \code{build\_member\_cache} from the
        \code{POST /BestTimes} responses and keyed by the memberId
        currently being processed.
\end{itemize}

% --------------------------- 4.5 ---------------------------------
\subsection{Interface view}

\paragraph{Command-line interface.}
\begin{itemize}
  \item \code{-o token[,token\ldots]} --- behaviour keywords and/or
        swimmer ids (see \S\ref{sec:why-ids}).
  \item \code{-e event[,event\ldots]} --- restrict to event codes;
        repeatable.
\end{itemize}

\paragraph{HTTP interface.}
Two endpoints on \code{https://times-api.usaswimming.org/swims/%
TimesSearch} are consumed (skipped under \code{offline}):
\begin{itemize}
  \item \code{POST /GetMembersForFilters} --- resolve a name to a
        \code{memberId} (body \code{\{"name":"\dots"\}}).
  \item \code{POST /BestTimes} --- best time per course for one
        stroke-and-distance pair (body
        \code{\{"memberId":"\dots","strokeAbbreviation":"\dots",%
        "distance":n\}}); each record carries \code{eventCode},
        \code{swimTime}, \code{swimDate}, \code{meetName}, and
        \code{timeStandard}.  Issued once per distinct pair.
\end{itemize}

Every request carries \code{Content-Type: application/json} plus the
anonymous data-hub auth headers \code{Usas-Sub-Id: Anonymous},
\code{AppName: DataHub}, and a run-time-generated \code{Device-Id}
(\S\ref{sec:why-anon}); there is no bearer token.

\paragraph{Database interface.}
Postgres is reached \emph{programmatically} through the
\texttt{libpq} client library (\req{REQ-I-12}) --- no
\texttt{psql} child process and no shell wrapper.  The
connection string is taken from
\texttt{SWIM\_TIMES\_PGCONNINFO} when set, otherwise the
default keyword string
\texttt{dbname=swim-times user=postgres host=localhost port=5432}.
A single \code{PGconn} is opened in \code{db\_open} and reused
for the lifetime of the run.  Two access patterns are needed:
\begin{itemize}
  \item Write (\code{store}): a prepared statement
        \texttt{INSERT INTO swim\_times (\dots) VALUES
        (\$1\dots\$7) ON CONFLICT (swimmer, event, swim\_time,
        swim\_date, meet) DO NOTHING}, executed once per row via
        \code{PQexecPrepared}.  Per-row inserts (rather than a
        single \texttt{COPY}) let \texttt{ON~CONFLICT~DO~NOTHING}
        absorb duplicates one at a time without rolling back the
        whole batch (\req{REQ-F-75}).
  \item Read (\code{offline}): one parameterised
        \texttt{SELECT \dots\ FROM swim\_times WHERE swimmer
        ILIKE \$1 AND event = \$2} per event via
        \code{PQexecParams}; tuples are read with
        \code{PQgetvalue}.
\end{itemize}

\paragraph{Standard output.}
Either a CSV stream (header line followed by data lines) or a
plain-text table (per-event heading, column header, separator,
data rows, blank line).

% --------------------------- 4.6 ---------------------------------
\subsection{Interaction view}

A typical invocation flows through the following sequence:

\begin{enumerate}
  \item \code{main} parses \code{-o} / \code{-e} via \code{getopt}.
  \item Unless \code{offline} is set, \code{main} initialises
        libcurl (\code{curl\_global\_init}).
  \item If either \code{store} or \code{offline} is set,
        \code{main} calls \code{db\_open} which opens a libpq
        connection (\code{PQconnectdb}) and prepares the
        idempotent \code{INSERT \dots\ ON~CONFLICT~DO~NOTHING}
        statement.  Failure under \code{offline} is fatal;
        failure under \code{store} only is logged and persistence
        is silently disabled.
  \item For each entry in \code{SWIMMERS} selected by
        \code{g\_sel} (all, if none), the dispatcher chooses the
        online or offline path:
        \begin{enumerate}
          \item \emph{Online.} \code{lookup\_member\_id} issues
                one POST and returns the \code{memberId}.
                \code{fetch\_times} builds that member's best-times
                cache once (one \code{POST /BestTimes} per distinct
                stroke-and-distance pair), then for each requested
                event (or all default events) selects the matching
                best time into a \code{TimeRow}, sorts, and emits to
                stdout, the DB, or both per the option mask.
          \item \emph{Offline.} \code{offline\_fetch} runs a
                parameterised \texttt{SELECT \dots\ FROM swim\_times}
                via \code{PQexecParams} for each event, copies
                each tuple into a \code{TimeRow} array, applies
                the date window, sorts, and emits.  No HTTP is
                performed; no child process is spawned.
        \end{enumerate}
  \item \code{main} calls \code{db\_close} (no-op when no
        connection was opened) and, unless \code{offline}, also
        calls \code{curl\_global\_cleanup} before exiting.
\end{enumerate}

% --------------------------- 4.7 ---------------------------------
\subsection{Algorithm view}

\paragraph{JSON scanning.}
\code{scan\_string} locates the next occurrence of the literal
\code{"key"}, advances past the colon and any whitespace, expects
a double-quote, and copies bytes until the matching closing quote
into the output buffer.  It advances an in/out position pointer so
that sequential scans walk forward through the response without
rescan, which is how \code{Parse BestTimes records} reads the fixed
\code{meetName}, \code{eventCode}, \code{swimDate}, \code{swimTime},
\code{timeStandard} field order of each record.  Every value the
program needs is a JSON string, so no numeric scanner is required.

\paragraph{Sort.}
\code{insertion\_sort} performs an in-place ascending sort of a
\code{TimeRow} array by \code{sort\_key}.  Chosen over heap or
quick sort because $n$ rarely exceeds 50 and insertion sort has
no recursion, no auxiliary memory, and trivial correctness.

\paragraph{Date reformat.}
\code{normalize\_date} converts the API's \code{"Mon DD, YYYY"} date
to ISO \code{YYYY-MM-DD} with a single \code{sscanf} and a
month-abbreviation lookup, bounding each field so the ten-byte result
always fits its buffer.  ISO order is lexical (so it sorts correctly),
contains no comma (so CSV needs no extra quoting), and matches the form
the offline path reconstructs with \code{to\_char}.

% --------------------------- 4.8 ---------------------------------
\subsection{Resource view}

\paragraph{Memory.}
Per-event peak memory is bounded by
$\code{MAX\_TIMES} \times \mathtt{sizeof}(\code{TimeRow})$
(approximately 70~KiB per event), plus the response buffer
(typically 30--80~KiB).  Buffers are freed at end of each
\code{fetch\_times} call; no long-lived heap state accumulates
across events.

\paragraph{Network.}
Each invocation issues at most one member-search POST plus one
best-times GET per selected swimmer (the best-times response is
cached and reused across that swimmer's events).  All connections are
short-lived; libcurl manages connection reuse internally where the
easy-interface allows.

\paragraph{Credentials.}
No credential is stored anywhere.  The anonymous data-hub headers
(\S\ref{sec:why-anon}) are assembled on each request and the
\code{Device-Id} is generated at run time by \code{device\_id()};
there is no static token to store or rotate.

% =================================================================
\section{Design Rationale}
% =================================================================

\subsection{Why CWEB?}

CWEB satisfies \req{REQ-A-04} (self-documentation) directly: the
typeset \code{swim-times.pdf} is the program's reference
documentation.  Equivalent C-with-comments would have required
maintaining design prose in a separate document, with the
inevitable drift between code and documentation that the
Parnas--Clements paper warns against.

\subsection{Why no JSON library?}

Adding \code{cJSON} or \code{jansson} would introduce a build
dependency and would obscure the small, focused parsing logic
that actually serves the program.  The response shape is
stable and the scanner has fewer than fifty lines of code; the
cost of a custom scanner is therefore lower than the cost of a
third-party dependency.  The scanner is deliberately format-agnostic:
it locates a quoted key and returns the value immediately following.
That contract held for the old Sisense JAQL arrays of
\code{\{"text":\dots\}}/\code{\{"data":\dots\}} cells and holds equally
for the current REST responses of named fields
(\code{"memberId"}, \code{"fullName"}, \code{"swimTime"}, \dots),
so the migration of \S\ref{sec:why-rest} reused it verbatim.

\subsection{Why the REST times API, not Sisense JAQL?}
\label{sec:why-rest}

The program was originally built against the Sisense JAQL analytics
API at \code{usaswimming.sisense.com} (a person-search cube and a
times Elasticube, each queried with a bearer token).  The USA~Swimming
data hub has since been rebuilt as a single-page app that talks to a
family of first-party REST services
(\code{times-api}, \code{person-api}, \code{meet-api},
\code{security-api}) behind an Azure API-management gateway; the
Sisense path is decommissioned for public use --- its embedded token
now returns \code{401 db\_unauthorized} and the current front-end
exposes no Sisense credential at all.  The fetch path was therefore
migrated to the public times service:

\begin{itemize}
  \item \code{POST /swims/TimesSearch/GetMembersForFilters} with body
        \code{\{"name":"\dots"\}} resolves a name to a member record
        carrying a \code{memberId} (an alphanumeric token such as
        \code{6CD35348E5824C}, replacing the old numeric
        \code{PersonKey});
  \item \code{POST /swims/TimesSearch/BestTimes} with body
        \code{\{"memberId":"\dots","strokeAbbreviation":"\dots",%
        "distance":n\}} returns that member's best time in each course
        (SCY and~LCM) for that stroke-and-distance pair, with swim
        date, meet name, and motivational standard
        (\S\ref{sec:why-bestonly}).
\end{itemize}

\noindent The endpoint contract was recovered from the site's
OpenAPI specification
(\code{https://times-api.usaswimming.org/swagger/v1/swagger.json})
and from its front-end bundle, then confirmed against the live service.

\subsection{Why anonymous headers, not a bearer token?}
\label{sec:why-anon}

The REST services do not use a bearer token for public data.  The
front-end's fetch wrapper authenticates every call with three custom
headers: \code{Usas-Sub-Id} (the caller's subject id, or the literal
\code{Anonymous} when not signed in), \code{AppName: DataHub}, and a
\code{Device-Id}.  The \code{Device-Id} is the only non-trivial part:
the server rejects a bare UUID (``\code{Invalid Device-Id format}''),
requiring the exact shape the browser produces --- base64 of
\code{"<platform> - <vendor> - <fingerprint> - <millis>"}, with the
first five base64 characters repeated after the fifteenth.
\code{device\_id()} reproduces that transformation over a per-run
fingerprint (process id plus clock); the server validates only the
format, not that the id is registered, so this is sufficient.  Because
the credential is generated at run time, there is nothing to embed and
nothing to rotate: the ``dynamically obtained valid token''
requirement is met by construction, and the failure mode of the old
static token (silent \code{401} after rotation) cannot recur.

\subsection{Why the \code{BestTimes} query, and how date, meet, and standard are recovered}
\label{sec:why-bestonly}

Two anonymous best-times feeds exist.  The lightweight
\code{GET /GetBestTimesForMember/\{memberId\}} returns one row per event
carrying only stroke, distance, course, and time --- no date, meet, or
standard --- and the program originally used it, leaving those three
\code{TimeRow} fields empty.  Requirements~\#1 and~\#2 of
\code{requirements6.txt} require each reported time to include its date,
its meet, and the \code{B}/\code{BB}/\code{A}/\code{AA}/\code{AAA}/%
\code{AAAA} standard attained, so the fetch path was moved to the
heavier \code{POST /swims/TimesSearch/BestTimes} query.  That query ---
the same one the data-hub ``Best~Times'' page issues --- is still
anonymous (it needs no signed-in subject), but its response records
carry \code{eventCode}, \code{swimTime}, \code{swimDate},
\code{meetName}, and \code{timeStandard}.  The still-richer per-swim
endpoints (\code{GetAllTimesForFilters}, \code{GetSwimmerMeets}, and
\code{person-api}'s \code{Person/search}) remain \code{403} for
anonymous callers and are not required for best-time reporting.

\code{BestTimes} keys on a single stroke-and-distance pair, so
\code{build\_member\_cache} issues one POST per distinct pair in the
effective event list (about twenty-two for the full thirty-one-event
roster, since each call returns both courses) and stores the parsed
rows in a per-member cache (\code{g\_cache}, keyed by
\code{g\_cache\_member}).  The per-event loop in \code{main} then reads
from that cache, so no event triggers a fresh network call.  Swim dates
arrive as \code{"Mon DD, YYYY"} and are normalised to ISO
\code{YYYY-MM-DD} by \code{normalize\_date} for lexical sorting,
comma-free CSV, and agreement with the offline read path.
\code{db\_insert\_row} still passes an empty field as SQL \code{NULL},
which now happens only for \code{standard} when a swim earned no
motivational standard; \code{schema.sql} accordingly keeps that column
nullable.

\subsection{Why id-based swimmer selection?}
\label{sec:why-ids}

The original design mapped each swimmer to a distinct bit
(\code{OPT\_STELLA}, \code{OPT\_KALEA}, \dots) in a single \code{int}
option word shared with the behaviour flags.  Extending the roster to
the full CAST squad would have overflowed that word.  Each
\code{Swimmer} now carries a unique string \code{id}; any \code{-o}
token that is not one of the four behaviour keywords (\code{fastest},
\code{csv}, \code{store}, \code{offline}) is treated as a swimmer id
and collected in \code{g\_sel}.  Selection is then a string match
against \code{SWIMMERS[].id}, with an empty selection meaning ``all''.
The roster can grow without touching the flag word, and the legacy
keywords (\code{stella}, \code{kalea}, \dots) continue to work because
they are simply the ids of the original entries.

\subsection{Why insertion sort?}

\req{D-5} addresses this directly.  The complexity argument is
not just about wall-clock time---it is about reading the code:
twelve lines of insertion sort can be verified by eye, whereas a
heap sort cannot.  For $n < 100$ the wall-clock difference is
imperceptible.

\subsection{Why one binary, not two?}

A separate person-search utility was considered but rejected.
The two queries (person search, times query) are always issued
together, and splitting them would force the caller to thread
the \code{PersonKey} through the shell, complicating every
invocation for no gain.

\subsection{Why \texttt{libpq} (programmatic), not \texttt{psql}?}

Earlier revisions reached Postgres by spawning a child
\texttt{psql} (via \texttt{pg.sh}) and either piping CSV into
\texttt{COPY \dots\ FROM~STDIN} or capturing pipe-separated
\texttt{SELECT} output.  Requirement~\#3 of
\texttt{requirements4.txt} reverses that decision: ``Access the
Postgres database programmatically, not through scripts.''  The
revised design therefore links \texttt{libpq} (Postgres's C
client library, autodetected via \texttt{pg\_config} in the
Makefile) and reaches the database through the standard
\code{PQ*} call surface --- \code{PQconnectdb},
\code{PQprepare}, \code{PQexecPrepared}, \code{PQexecParams}.
This (a)~eliminates the shell quoting and \texttt{popen}
\texttt{SIGPIPE} hazards of the previous design,
(b)~lets duplicate inserts be absorbed one row at a time by an
\texttt{ON~CONFLICT~DO~NOTHING} clause without aborting the
batch (the \texttt{COPY} path could not), and (c)~replaces the
brittle pipe-separated \texttt{SELECT} output with binary-safe
\code{PGresult} field access.  The cost is one new link-time
dependency on \code{libpq}; the SRS allowlist (\req{REQ-D-04})
is updated accordingly.

\subsection{Why per-row \texttt{INSERT \dots\ ON~CONFLICT}, not \texttt{COPY}?}

Requirement~\#4 of \texttt{requirements4.txt} demands that the
program ``check for duplicates before adding information \dots\
and recover appropriately if data already present.''  A bulk
\texttt{COPY \dots\ FROM~STDIN} cannot satisfy this: the first
row that violates the unique constraint aborts the entire
transaction, taking with it every subsequently-streamed row.
Switching to per-row \texttt{INSERT \dots\ ON CONFLICT
(swimmer, event, swim\_time, swim\_date, meet) DO NOTHING}
makes each insert independently idempotent: a duplicate is a
silent no-op, the run exits~0, nothing is written to standard
error, and any genuinely-new rows in the same invocation are
still persisted.  The cost is $N$ round trips instead of one
streamed transaction, but $N$ is at most a few hundred per
invocation and each insert completes in well under a
millisecond against the local container.

\subsection{Why a date-window field on \texttt{Swimmer}?}

\code{ledecky10}, \code{ledecky12}, and \code{ledecky14} originally
demanded age-bracket filtering, expressed as a \code{date\_min}/%
\code{date\_max} window on each \code{Swimmer} record rather than as a
separate command-line flag, so the per-swimmer policy lived in one
place and new historical comparisons were additive.  The fields are
retained on the record and in the \code{fetch\_times} signature for
compatibility, but they are not applied to the online fetch: the
\code{BestTimes} query (\S\ref{sec:why-bestonly}) returns a single
lifetime best per course, so the three Ledecky entries currently
resolve to the same best-times set.  The swim date is now reported for
each of those bests; restoring true age-bracket \emph{filtering} would
require the signed-in richer-times endpoint noted in
\S\ref{sec:why-bestonly}, whose per-swim history could be narrowed to a
date window.

\vfill
\begin{center}
  \small\textit{End of document CAST-SDD-001 Rev.\,2.0}
\end{center}

\end{document}
